\documentclass[conference]{IEEEtran}

\usepackage{cite}
\usepackage{amsmath,amssymb}
\usepackage{booktabs}
\usepackage{graphicx}
\usepackage{url}
\usepackage{microtype}

\title{Temporal Graph Diffusion Tracing for Cross-Modal Misinformation Factuality Degradation}

\author{
\IEEEauthorblockN{Arya Diwakar}
\IEEEauthorblockA{School of Computer Science\\and Engineering\\
Vellore Institute of Technology\\Vellore, India}
\and
\IEEEauthorblockN{Doravari Sakthi Prasad}
\IEEEauthorblockA{School of Computer Science\\and Engineering\\
Vellore Institute of Technology\\Vellore, India}
\and
\IEEEauthorblockN{Jeswanth G}
\IEEEauthorblockA{School of Computer Science\\and Engineering\\
Vellore Institute of Technology\\Vellore, India}
\and
\IEEEauthorblockN{Project Guide}
\IEEEauthorblockA{Dr. Jeevanantham A K\\
Professor (Higher Academic Grade)\\
School of Computer Science and Engineering\\
Vellore Institute of Technology, Vellore, India}}

\begin{document}
\maketitle

\begin{abstract}
Most propagation-based detectors treat a cascade as a fixed graph plus a root
post, and ignore what happens to the text representation as a claim moves from
one post to the next. We look at that edge-level change directly. Using the
official GossipCop and PolitiFact splits of UPFD, with frozen 768-dimensional
BERT features, we measure the distance between a child post's representation
and its parent's. A naive pooled comparison is dominated by the fact that real
cascades run deeper and wider than fake ones: on GossipCop the fake-to-real
ratio falls from 0.901 pooled to 0.977 once we weight by depth. The per-cascade
comparison (depth $\geq2$) gives 0.975, $p=0.048$ -- a small, borderline
effect -- while PolitiFact shows the same direction (0.969) but is not
significant ($p=0.522$). In five-fold cross-validation confined to the
GossipCop test partition, compact drift summaries reach 0.846 AUC against 0.792
for size alone ($\Delta=0.054$, bootstrap CI [0.040, 0.068]); this stays
inside the test split, so it is exploratory rather than an independent test.
Fake-labelled cascades also carry more near-duplicate sibling posts, and
sibling duplication correlates negatively with drift ($r=-0.310$), but the link
is observational and says nothing about coordination. Across five seeds the
drift-aware hybrid is not significantly better than an LSTM (accuracy
$p=0.832$, AUC $p=0.212$). The takeaway is that parent--child change is a real,
measurable signal -- not a proven detector improvement.
\end{abstract}

\begin{IEEEkeywords}
misinformation detection, propagation cascades, content drift, UPFD,
social-network analysis
\end{IEEEkeywords}

\section{Introduction}
A misinformation post is rarely shared verbatim. As it moves through a platform,
each repost rephrases, quotes, or trims the original, and the people who see it
add their own replies. The UPFD benchmark captures this as a tree: nodes are
posts, edges are reposts or replies, and every node carries a BERT embedding of
its text. What the benchmark does not make easy to see is the simplest question
we ask here -- when one post points to its parent, how much does the text
representation actually change?

We call that change \emph{parent--child content drift}. The guess behind it is
mundane: a cascade full of copied or templated posts should show less local
change than one with varied commentary. That is a tendency across cascades, not
a rule for any single post, and certainly not a way to read intent from
embedding similarity.

The title under which this work is submitted uses words we have had to pin down
narrowly. ``Temporal'' means topological depth in the tree, not wall-clock time.
``Cross-modal'' means the text-embedding view together with the
propagation-graph view -- not images, audio, or video. ``Factuality
degradation'' is not something we observe directly; we only compare embedding
change across labels that were already fact-checked. The point matters because a
representation drifting does not by itself show a post became less truthful.

Our contribution is narrower than the title suggests, and we keep it that way.
On UPFD \cite{dou2021upfd} we: (1) show why pooling all nodes hides a depth
confound and report the comparison after standardizing by depth; (2) compare
compact drift summaries with size and structure using cross-validation that
stays inside the test partition, labelled exploratory on purpose; (3) look at
sibling repetition as a correlate and report five-seed model runs without
claiming a win over a sequence baseline. We make no claim to state of the art,
to a causal mechanism, or to priority over earlier work.

Why look at the edge at all? Most public misinformation benchmarks ship a
cascade as a single labelled graph, so the natural unit for a detector is the
whole tree. The parent--child link is the one place where the benchmark already
tells us the direction of transmission, and it is exactly where a repost can
differ from what it replies to. We therefore treat the edge as the smallest
unit at which retransmission can be observed, while still reporting at the
cascade level so the statistics stay honest.

\section{Related Work}
Two lines of work sit behind this. The first builds detectors from propagation
structure. UPFD itself supplies the GossipCop and PolitiFact trees we use, with
node features ready for graph classification \cite{dou2021upfd}. BiGCN reads
propagation and dispersion directions in the tree \cite{bian2020bigcn}; earlier
recurrent models encode the order in which posts appear \cite{ma2016recurrent}.
These are prediction systems first.

The second line is the one our title loosely echoes. Fu \emph{et al.} use
``feature drift'' for how fake-news features shift across datasets or over time,
and study what that means for transferring a model \cite{fu2023drift}. Bai
\emph{et al.} fuse post meaning with the path and structure of propagation
using attention over tree nodes and edges \cite{bai2023tree}. GCN and GAT are
the standard aggregation tools \cite{kipf2017gcn,velickovic2018gat}, and BERT
supplies the contextual node features \cite{devlin2019bert}. None of these
measure the pairwise change between a post and its parent, which is the quantity
we study. We therefore borrow the word ``drift'' plainly; we are not presenting
it, or the broader use of text and structure, as novel.

The propagation structure has also been studied for early detection, where only
part of a cascade is observed at decision time. We return to that setting in the
partial-cascade experiments and explain there why it asks a different question
from the full-cascade comparison reported above.

\section{Data and Measurement}
\subsection{Benchmark and protocol}
We work with the UPFD release of GossipCop and PolitiFact, using the frozen
768-dimensional BERT features and the train/validation/test partitions exactly
as supplied \cite{dou2021upfd}. GossipCop gives us 1,092 training, 546
validation, and 3,826 test cascades; PolitiFact is far smaller at 62, 31, and
221. That tiny PolitiFact training set makes any neural-model number from it
unreliable, so we keep it only for the drift comparison, not for ranking
classifiers. The main models are trained on the fixed GossipCop split, with
the checkpoint chosen by validation accuracy, and each is rerun across five
seeds. Test cascades never enter checkpoint selection.

One example in this benchmark is a whole cascade, labelled real or fake, not an
individual post. Nodes are posts; edges are reposts or replies; the root sits
at depth zero. We take both the graph and the BERT features from the release
rather than rebuilding them from raw platform data. That keeps us comparable
with earlier UPFD work, but it also means we inherit whatever collection-time
or preprocessing choices the benchmark authors made. We are not reconstructing
a cascade's true timeline: depth is a topological order, not a guaranteed
timestamped sequence of posts.

For the primary neural comparison the train, validation, and test indices stay
fixed. Each seed sets the same initial seed in Python and the deep-learning
framework, trains on the training partition, and keeps the checkpoint with the
best validation accuracy. We then score that one checkpoint once on the test
partition for accuracy, F1, ROC-AUC, and average precision, and report the mean
and standard deviation over the five seeds. The paired tests line up results by
seed, so they describe variation across random initializations on this one
split -- not across different train/test partitions.

\subsection{Edge-level change and aggregation}
For a node $i$ with representation $x_i$ and parent $p(i)$, we define the two
edge distances
\begin{equation}
 d_i^{\mathrm{cos}}=1-\frac{x_i^\top x_{p(i)}}
 {\|x_i\|_2\|x_{p(i)}\|_2},\qquad
 d_i^{\mathrm{L2}}=\|x_i-x_{p(i)}\|_2 .
 \label{eq:drift}
\end{equation}
The root has no parent, so both are set to zero there. To summarize a cascade
we combine them into a single magnitude
$m_i=\sqrt{(d_i^{\mathrm{cos}})^2+(d_i^{\mathrm{L2}})^2}$. The L2 term scales
with the embedding, so we read this as a representation-space quantity rather
than a calibrated semantic distance; the cosine term rides along as a
scale-insensitive counterpart.

Pooling every node into one mean is misleading because it weights cascades by
their size and mixes depths together. We therefore report two cleaner views.
The first, a depth-standardized class mean, reweights each observed depth by
the pooled depth distribution so the two classes meet at the same depth mix.
The second is a cascade-level mean built only from nodes at depth $\geq2$, one
value per cascade, compared with an independent-samples test. Its $p$-value is
nominal -- we do not correct for the several descriptive comparisons we run.

Concretely, let $\bar m_{y,h}$ be the mean magnitude for class $y$ at hop
$h$, and let $q_h$ be the pooled share of nodes at that hop. The
depth-standardized mean is $\sum_h q_h\bar m_{y,h}$: it asks whether the class
means differ under a common depth composition, not whether one cause explains
the difference. The cascade-level version averages $m_i$ within each cascade
over nodes at hop two or deeper and then compares the per-cascade values.
Keeping hop one out stops the root's by-definition zero from artificially
pulling a cascade mean down, but it also means very shallow cascades drop out
of this estimand entirely; we report both class counts alongside the result.

\subsection{Compact feature-group diagnostics}
For each cascade we summarize the two content-change dimensions, along with
depth, normalized child count, and a root indicator, by their mean and standard
deviation, and we add log cascade size and maximum depth. A random forest of
200 trees is scored by stratified five-fold cross-validation over the official
test cascades, and a paired bootstrap compares its out-of-fold predictions on
the same cascades. Because the test partition does double duty here -- both for
the cross-validation and for the model -- these AUCs are exploratory looks
inside one split, not a clean estimate on unseen data.

The four content features are the mean and standard deviation of the two
parent--child distances. Structural features are the matching moments of
normalized hop depth, normalized child count, and the root indicator; size
features are log node count and maximum depth. We keep the groups apart on
purpose: if the concatenated set scores high, that should not be read as proof
that content change -- rather than cascade geometry -- is what separates the
classes. The forest is a diagnostic probe, not the main detector, and its
folds are stratified by cascade so that posts from one cascade never land in
two folds at once.

For the paired comparison we resample out-of-fold probabilities from the same
five folds by cascade and report the percentile interval on the content-only
minus size-only AUC gap. That interval captures resampling noise for this
partition and these folds; it does not fix the fact that we reused the test
split, and it is not a cross-platform uncertainty. No feature or threshold in
this retrospective probe is chosen on a separate validation set.

\subsection{Sibling repetition}
Within each cascade that has a sibling group, we count the share of sibling
pairs whose L2-normalized embeddings sit at cosine similarity of at least 0.95,
then relate that share to the cascade's drift. Similarity stands in for
repetition; it cannot tell us whether two posts were copied, coordinated, or
simply similar for some other reason.

Here the unit is the cascade, not each sibling pair. We fix the threshold at
0.95 cosine on L2-normalized embeddings and drop cascades with no sibling pair,
leaving 3,800 of the 3,826 GossipCop test cascades. The thresholded rate is an
interpretable proxy, and we use Pearson correlation against the cascade's
non-root mean drift. Moving the threshold shifts the absolute proportion -- the
project outputs include other cut-offs -- so we avoid reading the tiny
correlation $p$-value as a mechanism. With this many cases, even a modest
observational association can clear a significance bar.

\subsection{Neural models and training}
The comparison spans graph convolution (GCN), graph attention (GAT),
bidirectional propagation-graph convolution (BiGCN), a depth-ordered
bidirectional LSTM, DegradGNN, and HybridDegradLSTM. The LSTM baseline orders
node embeddings by propagation depth and runs them through a bidirectional
LSTM. DegradGNN projects the per-node cosine distance, L2 distance, depth,
child count, and root indicator into tree message passing. The hybrid joins a
sequence representation with a tree representation. These models let us ask two
different things: whether compact aggregate statistics track the label, and
whether an architecture that explicitly consumes drift helps. The hybrid test
is not a clean drift ablation, because it also changes the architecture and
adds auxiliary features.

All models take the 768-dimensional frozen BERT inputs, use a hidden dimension
of 128, two message-passing layers where relevant, and dropout 0.3. We train
with AdamW at learning rate $10^{-3}$ and weight decay $10^{-5}$, batch size
64, up to 20 epochs, and early stopping with patience 8, keeping the
checkpoint with the best validation accuracy. Each primary GossipCop model is
reported as mean and standard deviation over the same five seeds, and paired
$t$-tests compare seed-matched outcomes. With only five pairs those tests are
weak, so we treat them as rough checks rather than precise uncertainty, and we
deliberately avoid crowning a best seed.

In the hybrid, every 768-dimensional node vector is projected to 128 dimensions.
The depth-sorted sequence goes through a two-layer bidirectional LSTM with 128
units per direction, giving a 256-dimensional sequence summary. Alongside it,
the graph branch projects the five raw drift and structural values plus a
learned depth embedding, aggregates over two neighbor rounds, and mean-pools
the final 128-dimensional node states. A two-layer perceptron then classifies
the 384-dimensional concatenation. Root drift stays masked to zero while
rootness is carried as its own structural feature. This layout is why a
hybrid-versus-LSTM gap cannot be credited to the drift vector alone: the hybrid
also adds tree aggregation and rewires the fused representation.

The per-node drift fed to the network comes from the frozen input embeddings,
not from hidden states that shift during training. The edge distances we report
are therefore a fixed input-side diagnostic even though the classifier's own
representations are learned. Whether the signal holds up under a fine-tuned
encoder is left for future work. As shipped, the model sees the full benchmark
tree at prediction time unless we truncate it on purpose; it is not a streaming
classifier that updates after each new post.

\section{Results}
\subsection{Depth changes the interpretation of drift}
On GossipCop, real cascades average 75.8 nodes against 39.3 for fake-labelled
ones, and 80.4\% of real nodes sit at depth one versus 69.0\% of fake nodes.
Shallow nodes show larger changes on average, so when we pool everything the
fake-to-real drift ratio comes out 0.901 -- an overstatement of the
depth-controlled gap. Weight by depth and the ratio becomes 0.977, i.e. a 2.3\%
lower mean for fake-labelled cascades. Restrict the cascade-level view to depth
$\geq2$ and it is 0.975 ($p=0.048$, nominal; 1,910 real and 1,916 fake
cascades). The effect is small and the test is borderline, so caution is due.

The depth imbalance is not trivial: the GossipCop test set holds 219,989 nodes,
with 80.4\% of real-labelled but 69.0\% of fake-labelled nodes at hop one. At
hop one the mean magnitude is about 8.14 for real and 8.04 for fake, while over
nodes at hop two or deeper it is about 3.49 and 3.41. The depth mix therefore
shifts the pooled class means even though the within-depth gap is smaller. A
node-wise $t$-test across more than two hundred thousand observations would also
pretend posts from one cascade are independent, inflating precision; that is why
our main test aggregates to the cascade. The depth-standardized figure is
descriptive, and the cascade-level test is the inference we lean on.

The fake-to-real ratio sits below one at several common non-root depths -- 0.896
at hop two, 0.929 at hop three, 0.838 at hop four, 0.742 at hop five -- but the
pattern is not clean. At hop seven the estimate flips above one, on just 23 real
and 25 fake nodes. We do not read individual-hop tests as separate replications,
since later hops are sparse and we look at many of them. That scatter is another
reason to report one aggregate standardized estimate instead of claiming that
every fake cascade, or every depth, shows reduced change.

PolitiFact tells a similar story with weaker support: the depth-standardized
ratio is 0.911, yet the cascade-level ratio is 0.969 and not significant
($p=0.522$; 108 real and 113 fake cascades). The direction matches GossipCop,
but the cascade-level evidence does not hold up on this smaller set. This is
exactly why the pooled node statistic should not be read as a pure content
effect.

\begin{table}[t]
\caption{Depth-standardized and cascade-level drift comparisons. Ratios are
fake-labelled / real. Cascade-level test uses one mean per cascade at depth
$\geq2$.}
\label{tab:drift}
\centering
\footnotesize
\setlength{\tabcolsep}{2pt}
\begin{tabular}{lccc}
\toprule
Dataset & Depth-std. & Cascade-level & $p$ \\
\midrule
GossipCop & 0.977 & 0.975 & 0.048 \\
PolitiFact & 0.911 & 0.969 & 0.522 \\
\bottomrule
\end{tabular}
\end{table}

\subsection{Drift summaries and cascade size}
In the exploratory cross-validation on the official GossipCop test partition,
the content-only feature group reaches AUC 0.846 against 0.792 for size alone.
The paired out-of-fold gap is 0.054, with a percentile-bootstrap 95\% interval
of [0.040, 0.068]. Add content summaries to size and the AUC is 0.865; size
plus structure gives 0.850. The representation statistics do separate the
labels within this partition beyond size, but that is not the same as proving
they generalize: the features and the evaluation both live inside the official
test subset, so an untouched external set would be needed to confirm any
predictive value.

The full 12-feature diagnostic scores 0.879, and structural-only 0.853. That
fits the idea that the groups carry complementary information, yet the full-set
number does not say which group earns the credit, because the blocks correlate
and every estimate uses the same test partition. The bigger point is that size
alone already hits 0.792. So the fair reading is comparative: the edge-derived
summaries tell the labels apart better than these two size variables in this
diagnostic -- not that drift stands apart from every tree property, or that the
same lift will show up on fresh data.

We also tried nested groups: size plus content reaches 0.865 AUC, size plus
structure 0.850, a paired gap of about 0.015 with bootstrap interval
[0.005, 0.025]. Suggestive, but it carries the same within-partition caveat, so
we report it to make the comparison visible, not as a confirmed generalizable
edge for the content features.

\begin{table}[t]
\caption{Exploratory five-fold cross-validation within GossipCop's official
test partition. Results are not independent held-out estimates.}
\label{tab:diagnostic}
\centering
\footnotesize
\begin{tabular}{lcc}
\toprule
Feature group & AUC & Paired difference vs. size \\
\midrule
Size (2 features) & 0.792 & --- \\
Content drift (4 features) & 0.846 & $+0.054$ [0.040, 0.068] \\
Size + content drift & 0.865 & $+0.073$ \\
Size + structure & 0.850 & $+0.058$ \\
All 12 summaries & 0.879 & $+0.086$ \\
\bottomrule
\end{tabular}
\end{table}

\subsection{Repetition is associated with lower drift}
At cosine similarity $\geq0.95$, the mean sibling near-duplicate rate is 0.983
in fake-labelled cascades and 0.927 in real-labelled ones. Across the 3,800
usable cascades, duplication rate and drift correlate negatively
($r=-0.310$), and the sign holds within each label and after adjusting linearly
for log cascade size ($r=-0.312$). The broader within-cascade pairwise
similarity is almost the same across classes in the companion probe (0.9826
versus 0.9823), so the sibling result looks like localized repetition rather
than a uniform rise in all within-cascade similarity. Still, the correlation is
observational and stays inside the test partition; it points to no causal path
and no coordinated actors.

The threshold exposes how fragile this is. At 0.90 the near-duplicate rates are
already high in both classes (0.998 fake, 0.991 real); at 0.98 they are 0.880
and 0.822. That is why we anchor on the chosen 0.95 operationalization instead
of treating one arbitrary cut-off as natural. Mean pairwise similarity over all
posts in a cascade does not split the classes the same way, which hints that
local sibling repetition captures something different from global homogeneity.
The data do not test that causal story, though: the same similarity can come
from identical source wording, boilerplate, platform conventions, or plain
user copying.

\subsection{Model comparison and negative evidence}
Table~\ref{tab:models} lists the primary GossipCop runs. HybridDegradLSTM pairs
a depth-ordered bidirectional LSTM with a tree encoder that carries the drift
features; across five seeds its mean accuracy is 0.9564 and AUC 0.9874. The
LSTM baseline reaches 0.9573 and 0.9852. Paired tests over the matching seeds
find no significant gap in accuracy ($p=0.832$) or AUC ($p=0.212$), so we do
not claim the drift-aware hybrid beats the LSTM. It does beat DegradGNN, but
that comparison cannot pin the gain on the drift features, because the hybrid
also adds a sequence branch.

The paired accuracy difference, HybridDegradLSTM minus BaselineLSTM, averages
$-0.00094$ across seeds, while the AUC difference averages $+0.00224$. The
signs disagree by metric and neither test is significant. The standard
deviations in Table~\ref{tab:models} also show the runs are not rock-steady
across initializations: BiGCN's accuracy spread is 0.0238, far wider than the
hybrid's 0.0062. The ROC-AUC order may favor the hybrid, yet five seeds do not
suffice to call that order reliable or to credit it to drift. The honest
conclusion is the absence of clear evidence for an incremental hybrid benefit,
not a proof that the two approaches are equivalent.

\begin{table}[t]
\caption{Primary GossipCop test results, mean $\pm$ standard deviation over
five seeds. Fixed official split; validation-accuracy checkpoint selection.}
\label{tab:models}
\centering
\footnotesize
\setlength{\tabcolsep}{3.3pt}
\begin{tabular}{lcc}
\toprule
Model & Accuracy & ROC-AUC \\
\midrule
GCN & $0.7841\pm0.0110$ & $0.8694\pm0.0080$ \\
GAT & $0.9165\pm0.0046$ & $0.9709\pm0.0019$ \\
BiGCN & $0.9303\pm0.0238$ & $0.9772\pm0.0097$ \\
DegradGNN & $0.9294\pm0.0029$ & $0.9771\pm0.0006$ \\
BaselineLSTM & $0.9573\pm0.0034$ & $0.9852\pm0.0019$ \\
HybridDegradLSTM & $0.9564\pm0.0062$ & $0.9874\pm0.0017$ \\
\bottomrule
\end{tabular}
\end{table}

An earlier single-run component ablation (seed 1) found that stripping per-node
drift and hop features actually lifted DegradGNN accuracy from 0.9247 to 0.9378.
That is a weak, single-run signal, but it is enough to warn against treating
drift as a node feature that always helps. A separate three-seed partial-cascade
study showed all compared models losing substantial accuracy at 10\% of the
cascade, and root-only truncations carry zero drift by construction. Together
they argue for measuring drift in aggregate while leaving a robust early-warning
design an open problem.

For clarity, that ablation compares full DegradGNN (accuracy 0.9247, ROC-AUC
0.9705) with the same run minus drift (0.9378, 0.9829), minus depth (0.9326,
0.9787), and minus both auxiliary groups (0.9365, 0.9796). A no-tree MLP scores
even higher (0.9634, 0.9866), which tells us this one run does not isolate a
message-passing benefit on the chosen split. None of these are averaged over
seeds, so they should not be lined up numerically against the five-seed primary
table as if they carried equal weight. We include them as a counter to the
tempting but unsupported claim that dropping drift into every node must improve
classification.

In the previously checked three-seed partial-cascade test, models trained on
full cascades are then scored after keeping only the shallowest fraction of each
test cascade. At 10\% mean accuracy is 0.787 for the LSTM, 0.814 for DegradGNN,
and 0.808 for the hybrid; by 25\% all three sit near 0.91. The 10\% gap is
descriptive only and was not significant with three seeds. For the compact
feature probe, content-only AUC at 10\% is 0.803 with all truncated cascades
included but drops to 0.718 once root-only cases are removed. About 25.7\% of
truncated test cascades contain only the root, where parent--child drift is
zero by definition. That boundary means partial-cascade drift scores mix a real
content signal with the mere presence of at least one edge, so they must not be
sold as reliable early detection on their own.

These early numbers answer a different question from the five-seed full-cascade
evaluation: the models train on complete cascades and are then tested on
cascades cut to a fraction of their nodes. The setup measures no latency from
event onset, because timestamps and an online intervention policy are absent,
and it trains no separate classifier per observation fraction. The 10\% figures
are therefore a stress test of missing input, not evidence of operational early
detection. A real deployment study would need timestamp-faithful truncation,
fixed observation windows, calibration, and an error-cost analysis.

\section{Discussion and Limitations}
On GossipCop the descriptive finding is modest: after depth is controlled,
representation change along parent--child edges still differs by label, but the
effect is small and does not reach significance on PolitiFact. The sibling
result is one possible reading -- more repeated sibling content, lower drift --
yet it cannot tell copying apart from shared source text, platform conventions,
or other confounders, and we have no coordination labels to check against.

Several more limits are worth stating plainly. Both datasets lean on benchmark
labels and platforms, so label noise, collection practice, and
class-conditional tree construction can all leak into the results. The L2 term
is sensitive to embedding scale. The compact-feature probe reuses the official
test partition for cross-validation and must stay exploratory, not a clean
generalization test. PolitiFact is small and its neural numbers are unstable.
And the five-seed primary comparison lives on one benchmark split while the
ablation is single-run and the early-detection work is three-seed. Together
these rule out any claim of broad robustness or operational readiness.

The benchmark is also observational and historically bounded. GossipCop and
PolitiFact were built around fact-checking labels, not drawn as random samples
of social media, so their cascades may differ in coverage, audience, topic,
platform behavior, and moderation history. A label-linked embedding pattern
could therefore reflect those selection choices as much as message
retransmission, and the present comparison cannot pull the two apart. A
replication worth its name would use independently collected data, keep
timestamps and parent links, document how labels were adjudicated, and test
across time and topic rather than within a single random split.

Where the cascade is the right scientific unit, we keep the inference at that
level, so a few very large trees cannot dominate the significance test just by
holding many posts. That still leaves confounding by source, topic, language,
or collection process untouched, and the cross-dataset comparison is not a
pooled test: PolitiFact's small test count makes it far less precise. We
therefore report its direction and its nonsignificant cascade-level $p$-value
rather than fusing the two datasets into one headline number.

All primary classifier scores come after checkpoint selection on the fixed
validation split; we do not pick the outcome among test seeds. The diagnostic
random forest is a different animal and reuses the official test partition in
five-fold cross-validation, meant to compare feature summaries after the fact
and to expose the cascade-size confound. A confirmatory study should lock the
feature construction and hyperparameters on training and validation first, then
evaluate exactly once on fresh data. Without that separation a diagnostic AUC
can too easily be mistaken for a leaderboard result.

The reproducibility bundle carries the analysis scripts, the JSON summaries,
and the seed-level primary results. The six main classifier runs finished all
five seeds; an optional DriftPool rerun stopped at four and is left out of the
complete-seed table. Early-detection curves stay three-seed and the component
ablation stays single-run, as noted. Stating these scopes stops a partial run
from quietly looking comparable to the finished primary experiment. We do not
promise bitwise determinism across hardware backends, so the scripts and saved
outputs -- not an assertion of identical bits -- define what is reproducible.

The terminology overlap deserves a clear line. Earlier work uses feature drift
for change across datasets or time; we measure change within a cascade between
neighboring posts. That is a difference in operationalization, not a claim that
no one has studied semantic change in propagation before. A systematic
literature pass and an independent reproduction would be needed before any
priority claim.

The practical upshot is accordingly narrow. A depth-aware drift summary could
be offered to a human analyst as one extra signal, but the current results do
not justify automating accusations, removing content, or flagging accounts as
coordinated. Similarity can spring from ordinary quotation and shared sources,
and the class-level effect is small. In a human-review workflow, a candidate
use would have to show the summary improves calibrated decisions at a fixed
false-positive rate and that the gain survives across topics and collection
periods -- a decision analysis we have not run.

One further caution concerns the labels themselves. Both benchmarks inherit
fact-checking verdicts and treat them as if they were ground truth, yet those
verdicts are judgments made by particular outlets at particular times, and they
say nothing about whether a given post is harmful, illegal, or in need of any
action. A drift summary that separates real from fake in aggregate therefore
speaks only to measured representation change; it carries no recommendation
about what a platform or moderator should do with any individual post. We keep
the analysis at that descriptive level and avoid any policy claim.

A stronger follow-up would keep discovery and confirmation apart. Feature
definitions, depth weighting, and the sibling-similarity threshold could be
fixed up front on a development corpus, with a second independently collected
corpus held back for a single evaluation. Timestamped parent links would allow
comparison at matched elapsed time as well as matched depth, since equal depth
is not equal age. Annotating a subset by hand could separate verbatim copying,
quotation, paraphrase, and unrelated siblings; if coordination is the question,
that annotation must rest on independent behavioral evidence, not on similarity
inferred from embeddings. Finally, model comparisons should freeze the sequence
and graph architecture while varying only the drift representation, draw on
multiple random splits or external datasets, and report calibration and
class-conditional error alongside rank metrics. Those steps would turn today's
benchmark observation into a testable story about how content transformation
travels with propagation.

\section{Conclusion}
What looked like a large gap between fake and real cascades on GossipCop
narrows to a small one once depth is taken seriously: the pooled 0.901
fake-to-real ratio becomes 0.977 after depth weighting, and 0.975 at the
cascade level ($p=0.048$, nominal). PolitiFact points the same way (0.969) but
is not significant ($p=0.522$). The drift summaries do carry information past
cascade size in our within-partition cross-validation (0.846 versus 0.792 AUC),
and sibling repetition tracks lower drift ($r=-0.310$), but both findings are
exploratory and observational. Across five seeds, the drift-aware hybrid does
not beat the LSTM (accuracy $p=0.832$, AUC $p=0.212$). The honest reading is
that parent--child change is real and measurable, not that it is a detector
waiting to be deployed. Whether it generalizes is a question for independent,
timestamped data and a preregistered plan.

\begin{thebibliography}{00}
\bibitem{dou2021upfd}
Y. Dou, K. Shu, C. Xia, P. S. Yu, and L. Sun, ``User preference-aware fake
news detection,'' in \emph{Proc. 44th Int. ACM SIGIR Conf. Research and
Development in Information Retrieval}, 2021, pp. 2051--2055.

\bibitem{bian2020bigcn}
T. Bian, X. Xiao, T. Xu, P. Zhao, W. Huang, Y. Rong, and J. Huang, ``Rumor
detection on social media with bi-directional graph convolutional networks,''
in \emph{Proc. AAAI Conf. Artificial Intelligence}, vol. 34, no. 1, 2020,
pp. 549--556.

\bibitem{ma2016recurrent}
J. Ma, W. Gao, P. Mitra, S. Kwon, B. J. Jansen, K.-F. Wong, and M.-Y. Cha,
``Detecting rumors from microblogs with recurrent neural networks,'' in
\emph{Proc. 25th Int. Joint Conf. Artificial Intelligence}, 2016,
pp. 3818--3824.

\bibitem{fu2023drift}
C. Fu \emph{et al.}, ``Feature drift in fake news detection: An interpretable
analysis,'' \emph{Applied Sciences}, vol. 13, no. 1, Art. no. 592, 2023,
doi: 10.3390/app13010592.

\bibitem{bai2023tree}
L. Bai, X. Han, and C. Jia, ``A rumor detection model incorporating
propagation path contextual semantics and user information,''
\emph{Neural Processing Letters}, vol. 55, 2023, doi:
10.1007/s11063-023-11229-w.

\bibitem{kipf2017gcn}
T. N. Kipf and M. Welling, ``Semi-supervised classification with graph
convolutional networks,'' in \emph{Proc. Int. Conf. Learning Representations},
2017.

\bibitem{velickovic2018gat}
P. Veli\v{c}kovi\'c, G. Cucurull, A. Casanova, A. Romero, P. Li\`o, and
Y. Bengio, ``Graph attention networks,'' in \emph{Proc. Int. Conf. Learning
Representations}, 2018.

\bibitem{devlin2019bert}
J. Devlin, M.-W. Chang, K. Lee, and K. Toutanova, ``BERT: Pre-training of
deep bidirectional transformers for language understanding,'' in
\emph{Proc. NAACL-HLT}, 2019, pp. 4171--4186.
\end{thebibliography}

\end{document}
